% P8 (fraud side-probe) standalone abstract
\begin{abstract}
As a deliberate cross-domain probe of the \ours{} measurement discipline
---\emph{measure capability, do not trust the label}---this side study steps
outside GRPO to stress-test the same principle in a setting with a strong
non-LLM baseline. Where do large language models (LLMs) actually add value in
credit-card fraud detection when a tuned gradient-boosted tree already scores well? We run a
head-to-head on a custom synthetic single-configuration fraud dataset (50{,}000
transactions, $\approx$1.4\% realized fraud rate). The empirical results are negative for LLM scoring:
XGBoost reaches test AUC $0.7955$ on the $10{,}000$-row held-out split, while
a Qwen3.5-4B SFT row-serialization arm reaches accuracy $0.792$ but AUC
$0.48268$ on a 500-row positive-enriched held-out evaluation. (The two AUCs sit on
different held-out splits and are not a like-for-like ranking comparison; the
load-bearing point is only that the LLM scorer is at chance-level ranking, not that
it is worse than the tree by this exact margin.) The gradient-boosted tree keeps the real-time scorer seat for five compounding reasons:
artifact-backed AUC, latency, cost, calibration, and prompt-injection exposure.
But framing the comparison as
``LLM \emph{versus} XGBoost'' asks the wrong question. We identify four
capability gaps where the LLM contributes something the tree cannot \emph{natively}
represent or generate (even if hand-engineered features could be fed to a tree):
(i) document and image fraud via vision--language models, whose
raw pixel-and-text inputs lie outside any tree's feature space; (ii) compliance
narration---Suspicious Activity Report narratives and adverse-action
letters---a regulated text-generation category (FinCEN; ECOA/Regulation~B);
(iii) cold-start triage of novel fraud typologies before labels exist; and
(iv) agentic investigation of the alert queue, at roughly $85\times$ lower cost
than a human analyst by our estimates. We distill these findings into a
hybrid architecture in which the LLM serves as \emph{sensor} (feature extractor
over unstructured evidence) and \emph{scribe} (regulatory narrator), while
XGBoost remains the \emph{scorer}, with post-score agentic triage on the alert
queue. The study is the side-probe of a reproducibility program whose thesis is
``measure capability, don't trust the label'': applied here, the label under
test is ``AI.''
\end{abstract}
